\section{Background and Related Work}
\label{sec:background}
\label{sec:related_work}

A configurable system exposes options whose valid combinations form a configuration space $C$. For a fixed workload and non-functional property, each
configuration $c\in C$ has a measured value $y_c$ such as runtime, output size,
or energy. Sampling-based performance prediction trains a regression model
$\hat y$ from a measured subset $S\subset C$ and evaluates prediction error,
typically mean relative error
($\MRE=(1/|T|)\sum_{c\in T}|y_c-\hat y_c|/|y_c|$, lower is better), on
held-out configurations $T$~\cite{siegmund2015FSE,kaltenecker2019,pereira2020,julianasurvey}.

\paragraph*{Sampling strategies}
The configurable-systems literature offers several domain-agnostic families.
Random sampling draws configurations uniformly. Solver-based sampling asks
SAT/SMT solvers (Sat4j, Z3) for valid configurations, and its randomized variant
perturbs solver order for diversity \cite{DBLP:journals/tse/HenardPPKHT14}.
$t$-wise sampling covers option-value combinations
\cite{Kuhn2004,Cohen2008,Johansen2012}. Distance-based and diversified
distance-based sampling spread the sample in feature space \cite{kaltenecker2019}.
When artifacts provide measured pools rather than feature-model handles, we use
pool-based analogues: \emph{balanced} keeps option frequencies close to the pool,
\emph{coverage} matches marginal option coverage, \emph{diversity} spreads by
Hamming distance, \emph{uncertainty} adds high-residual configurations,
\emph{pairwise} covers option-value pairs, and \emph{pilot-cost} learns a small
pilot cost model before choosing configurations predicted inexpensive.

\paragraph*{How prior work counts effort.}
Sampling-based performance prediction has studied random, solver-based,
coverage-based, distance-based, progressive, and transfer-learning strategies
\cite{guo2013ASE,sarkar2015,siegmund2015FSE,nair2018c,kaltenecker2019,julianasurvey},
and almost all of them equate measurement effort with sample size. Sarkar
\etal, the closest cost-aware precedent, define a composite objective with
training, testing, model, and error cost, but collapse acquisition effort to a
term proportional to $n$ under unit per-configuration cost \cite{sarkar2015}.
Westermann \etal discuss the number of measurement points needed for a target
error \cite{westermann2012}. Guo \etal vary measurement budgets ($N$, $2N$,
$3N$, and a larger reference $M$) \cite{guo2013ASE}. Siegmund \etal report
Pareto trade-offs between accuracy and number of measurements
\cite{siegmund2015FSE}. Kaltenecker \etal compare strategies at common
sample sizes \cite{kaltenecker2019}. Siegmund \etal also report elapsed hours
for some heuristics \cite{siegmund2012a}, but not as a test of whether $n$ is a
valid cost proxy. To our knowledge the proxy itself has not been evaluated in
this literature: our contribution is not a new sampler but a validity study of
that effort axis.
